\Needspace{10\baselineskip}
\section{Phân tích yêu cầu hệ thống}
\label{sec:yeu-cau}
\subsection{Các tình huống sử dụng chính}
Các yêu cầu được xác định từ cách người dùng trao đổi với chatbot:
\begin{itemize}[leftmargin=*,itemsep=2pt,topsep=3pt]
\item Khi câu hỏi đã có đủ thông tin và căn cứ, người dùng nhận hướng dẫn trực
 tiếp, kèm nguồn hỗ trợ và bước tiếp theo.
\item Khi thiếu thông tin hoặc cách diễn đạt còn mơ hồ, người dùng được hỏi
 đúng dữ kiện có ảnh hưởng đến hướng dẫn.
\item Khi bổ sung, sửa hoặc rút lại một thông tin, người dùng nhận phản hồi
 theo tình huống mới và biết phần hướng dẫn nào đã thay đổi.
\item Khi nguồn chưa đủ, còn mâu thuẫn hoặc yêu cầu ngoài phạm vi, người dùng
 được thông báo phần chưa thể kết luận và đầu mối cần trao đổi.
\end{itemize}
Những tình huống này không tạo thành một kịch bản bắt buộc cho mọi câu hỏi.
Câu tra cứu đơn giản đã đủ căn cứ có thể được trả lời ngay.

\subsection{Yêu cầu chức năng}
Bảng~\ref{tab:chuc-nang} chuẩn hóa các hành vi cần bảo đảm. Đây là cơ sở để
lựa chọn thành phần thiết kế và xây dựng tình huống kiểm thử.

\begin{table}[H]
\centering\small
\begin{tabularx}{\textwidth}{|P{13mm}|P{33mm}|Y|}
\hline
\textbf{Mã} & \textbf{Yêu cầu} & \textbf{Hành vi cần bảo đảm} \\\hline
F01 & Hỏi bổ sung có mục đích & Chỉ hỏi thông tin thiếu có thể đổi hướng dẫn; không hỏi lại dữ kiện đã rõ. \\ \hline
F02 & Đối chiếu điều kiện & Dùng đúng nguồn cho từng phần của cùng tình huống; không ghép thêm miền không liên quan. \\ \hline
F03 & Cập nhật khi dữ kiện đổi & Chọn lại thẻ và nguồn, kiểm kết luận phụ thuộc; đánh dấu kết luận đã được thay thế. \\ \hline
F04 & Phân biệt mức kết luận & Trả đáp ứng theo khai báo, không đáp ứng theo khai báo hoặc chưa đủ căn cứ cho từng điều kiện. \\ \hline
F05 & Giải thích kiểm tra được & Chỉ rõ dữ kiện, điều kiện, đoạn nguồn, phiên bản hoặc mốc áp dụng và bước tiếp theo. \\
\hline
\end{tabularx}
\caption{Yêu cầu chức năng F01 đến F05}
\label{tab:chuc-nang}
\end{table}

Ba mức kết luận ở F04 được gắn với từng điều kiện đang xét. Việc xác định
một nguồn có liên quan không được diễn đạt thành xác nhận toàn bộ hồ sơ hợp
lệ hoặc một khoản chi chắc chắn được thanh toán.

\subsection{Yêu cầu phi chức năng}
Hệ thống cần hỗ trợ câu hỏi tiếng Việt, cho phép người dùng hỏi tiếp và sửa
thông tin bằng hội thoại. Cách trình bày phải giúp người đọc nhận ra hướng
dẫn chính, phần còn thiếu và căn cứ tương ứng mà không phải tự tìm trong
một danh sách tài liệu chung.

Thông tin đang dùng và quan hệ giữa dữ kiện, điều kiện, kết luận cần truy
vết được. Điều này phục vụ kiểm tra câu trả lời và cập nhật nội dung khi
nguồn thay đổi. Các chỉ tiêu về độ chính xác và thời gian phản hồi sẽ được
xác lập, đo và phân tích trong phần thực nghiệm.

Nếu sử dụng tài liệu nội bộ, quyền truy cập phải được kiểm soát trước truy
xuất và trước hiển thị. Dữ kiện nhạy cảm chỉ được thu thập khi cần cho tình
huống; thời hạn lưu và phạm vi xem lịch sử phải phù hợp với quyền sử dụng
nguồn và thông tin người dùng.

\subsection{Yêu cầu dữ liệu}
Văn bản gốc, tài liệu phân tích, nội dung đã chuẩn hóa và tình huống giả định
cần được phân biệt rõ. Nguồn dùng để trả lời phải được kiểm tra về xuất xứ,
độ đầy đủ, phiên bản, hiệu lực, đối tượng và quyền sử dụng. Tài liệu định
hướng hỗ trợ tổ chức nội dung, không thay thế căn cứ chính thức.

Điều kiện tư vấn phải dẫn được tới đoạn nguồn đã kiểm tra và có đầu mối
xác nhận nội dung. Bộ ca kiểm thử cần chỉ rõ dữ kiện đầu vào, nguồn áp dụng,
đáp án mong đợi và điểm cần hỏi bổ sung hoặc dừng. Mục~\ref{sec:kien-truc}
trình bày cấu trúc thông tin và cách tổ chức tri thức để đáp ứng các yêu
cầu trên.
